\section{Truppengattungen}
	In diesem Abschnitt werden die einzelnen Strukturelemente und die Organisation des \ac{TTT} während einer typischen Mission erklärt. Dabei wird auf das Buddyteam, den Trupp, den Zug und die gesamte Struktur der Kräfte eingegangen.\\
	Die Struktur lehnt sich dabei an das System der Bundeswehr sowie der Amerikanischen Streitkräfte an und wird im \ac{TTT} als 6 + 2 System so gelebt.\\
	Ganz unten in der Hierarchie der Strukturelemente steht das Buddyteam, welches aus zwei Soldaten besteht. Sie bilden die kleinste, elementarste Einheit. Zwei oder mehr Buddyteams werden durch einen Truppführer koordiniert und bilden so einen Trupp. Zwei Trupps bilden wiederum einen Zug, der durch die Zugführung geleitet wird. Alle Züge und Elemente werden von den Anweisung seitens der Operativen Planung (\ac{OPL}) geleitet.  

\subsection{Operationsleitung}
\subsection{Zugführung}
\includegraphics[width=20mm]{./Grafiken/Abschnitt/TrBraun}\quad\includegraphics[width=20mm]{./Grafiken/Abschnitt/TrGruen}\\
Zugführung -- Trupp Grün und Trupp Braun -- befehligt Fallschirmjägertrupp Schwarz und Pioniertrupp Blau, übernimmt im Notfall medizinische Erstversorgung und koordiniert \ac{MedEvac}, \ac{CAS} und Logistik wenn benötigt. 
\subsection{Infanterie}
\subsection{Spezialkräfte}
\includegraphics[width=20mm]{./Grafiken/Abschnitt/TrGold}\\
Gold -- Fernspäher -- führt Aufklärung mit Drohnen und Scharfschützenteams und bekämpft Einzelziele mit hoher Priorität.\\\\
\includegraphics[width=20mm]{./Grafiken/Abschnitt/TrGrau}\\
Grau -- Kommando Spezialkräfte -- Spezialkräfte zur gesonderten Verwendung, welche parallel zu den regulären Kräften operieren und deren Vorgehen durch Aufklärung, Sabotage und gezielte Bekämpfung von Einzelzielen unterstützt.

\subsection{Logistik}
\includegraphics[width=20mm]{./Grafiken/Abschnitt/TrSilber}\\
Silber -- Bussard -- Stellt Fahrzeuge und Personal bereit, mit denen Transport und Logistik durchgeführt werden.
\subsection{MedEvac}
\includegraphics[width=20mm]{./Grafiken/Abschnitt/TrWeiss}\\
Weiß -- \acf{MedEvac} -- Unterstützt Operationen mit Versorgungs- und Transportkapazität für Verwundete.
\subsection{Close Air Support}
\subsection{Mechanisierte Infanterie}
\subsection{Kampfpanzer}
\subsection{Artillerie}